% SHARD-ID: RC-02I-DEFINITIONS-CCM-TENSOR
% SHARD-TITLE: Definitions, Continued: Transfer Histories, Weil Kernels, and CCM Quotients
% SHARD-SUMMARY: Distinguishes the physical MPS auxiliary space, doubled transfer space, operator feature space, and cyclic moment space.
% SHARD-SUMMARY: Fixes the metric, conjugation, window and boundary conventions for the tensor-network reconstruction of finite and continuous spectral data.
% SHARD-KEYWORDS: CCM, tensor network, Gram matrix, clock, feature map, moment space, quotient, kernel
% SHARD-NOTES: none
% SHARD-SCRIPTS: none

\section{Definitions, Continued: Transfer Histories and CCM Quotients}
\label{sec:definitions-ccm-tensor}

These definitions supplement \cref{def:transfer-matrix,def:weil-form};
they do not identify the physical and reconstructed auxiliary spaces.
Inner products are conjugate-linear in the first argument.

\begin{definition}[The four spaces of transfer reconstruction]
\label{def:ccm-tn-spaces}
Let $\TNaux$ be the physical MPS auxiliary space. The transfer matrix of
\cref{def:transfer-matrix} acts on
$\TNspace=\TNaux\otimes\overline{\TNaux}\simeq\operatorname{End}(\TNaux)$.
A retained invariant space $\TNret$ is specified separately; when spectral
subtraction removes partial multiplicities, only a retained spectral multiset
is specified and an invariant restriction is not thereby defined.
Write $E$ for the rescaled retained operator when such an operator is given,
and $d=\dim\TNret\geq1$; empty retained spaces are excluded in these
history-state constructions. The feature space is
$\operatorname{End}(\TNret)$; the reconstructed moment space $\TNmoment$
will be a quotient of polynomials, not a synonym for $\TNaux$ or $\TNret$.
For an edge lift, $\TNspace$ includes the directed-edge or letter register.
\end{definition}

\begin{definition}[Metric adjoint and transfer-feature map]
\label{def:ccm-tn-feature}
Given $\TNmetric>0$ on $\TNret$, put
$A^{\sharp}=\TNmetric^{-1}A^*\TNmetric$ and
$\langle A,B\rangle_{\mathrm{HS},G}=\Tr(A^\sharp B)$.
Equivalently this is the ordinary Hilbert--Schmidt product of
$\widetilde A=\TNmetric^{1/2}A\TNmetric^{-1/2}$ and $\widetilde B$.
For a window of $K\geq1$ lengths $j=0,\ldots,K-1$, the length register is
$\mathbb C^K$ and the feature map is
\[
 \TNfeature c=\sum_{j=0}^{K-1}c_jE^j=p_c(E),\qquad
 p_c(z)=\sum_{j=0}^{K-1}c_jz^j.
\]
Vectorisation is an isometry from the ordinary Hilbert--Schmidt operator
space to two open indices. A bar on this vector belongs to the conjugate
feature Hilbert space; this convention fixes the transpose in a clock
partial trace.
\end{definition}

\begin{definition}[Moment matrix, cyclic algebra, and shift]
\label{def:ccm-tn-moments}
For the retained rescaled trace sequence of
\cref{def:rescaled-trace-sequence}, write $t_k=\Tr E^k$ for $k\geq0$
and $t_{-k}=\overline{t_k}$. Its window matrix is
\[
  (\TNGram)_{jk}=t_{k-j},\qquad 0\leq j,k<K.
\]
Thus $c^*\TNGram c=\Wform(c)$ with the coefficient convention of
\cref{def:weil-form}. When the spectral support is a finite set
$\{z_1,\ldots,z_r\}$ on the unit circle with positive weights $m_a$,
put $p_*(z)=\prod_{a=1}^r(z-z_a)$ and
\[
 \TNmoment=\mathbb C[z]/(p_*),\qquad
 \langle[p],[q]\rangle=\sum_a m_a\overline{p(z_a)}q(z_a).
\]
The moment shift is multiplication by $z$ on this quotient.
The distinct-support count $r$ and the total weight $\sum_a m_a$ are
different data. An unweighted trace of the moment shift counts support once.
\end{definition}

\begin{definition}[Transfer history and propagation Hamiltonian]
\label{def:ccm-tn-history}
In \cref{def:ccm-tn-feature}, let
$v_j=\operatorname{vec}(\TNmetric^{1/2}E^j\TNmetric^{-1/2})$ and
\[
 \TNhistory=\sum_{j=0}^{K-1}|j\rangle\otimes\overline{v_j}.
\]
When $E^*\TNmetric E=\TNmetric$, let $V$ act on the conjugate feature
space as the conjugate of left multiplication by
$\TNmetric^{1/2}E\TNmetric^{-1/2}$. For $0\leq j<K-1$ put
$L_j=\langle j+1|\otimes I-\langle j|\otimes V$ and
$H_{\mathrm{prop}}=\sum_jL_j^*L_j$.
The clock Hamiltonian is $H_{\mathrm{clock}}=\TNGram$ when the latter
is positive, and $P_{\ker}$ denotes its orthogonal kernel projector.
These operators act on a length register, or a length register with a
feature space; no locality in the physical sites of the original MPS is
part of these definitions.
\end{definition}

\begin{definition}[CCM finite window, shifted metric, and boundary correction]
\label{def:ccm-tn-loewner}
On the Fourier coefficient space $\mathbb C^{2N+1}$, indexed by
$-N,\ldots,N$, let $D=\operatorname{diag}(j)$,
$\eta=(1,\ldots,1)^T$, and $Jv=(v_{-j})_j$.
A CCM matrix $Q_N$ is real symmetric, with
\[
 (Q_N)_{jj}=a_j,\qquad (Q_N)_{jk}=\frac{b_j-b_k}{j-k}\ (j\ne k),
 \quad a_{-j}=a_j,\quad b_{-j}=-b_j,\quad \TNresponse=(b_j)_j.
\]
Set $\epsilon=\lambda_{\min}(Q_N)$ and $T=Q_N-\epsilon I$.
The even-simple hypothesis means that this least eigenvalue is simple
and its eigenvector satisfies $J\TNnull=\TNnull$. When $\eta^*\TNnull=1$,
the corrected operator is $D'=D-|D\TNnull\rangle\langle\eta|$.
The proposed quotient operator $D''$ acts on
$\mathbb C^{2N+1}/\mathbb C\TNnull$ with metric $T$; its well-definedness
and self-adjointness are assertions to prove, not part of the definition.
In the interval $[0,L]$ the modes are
$u_j(x)=L^{-1/2}e^{2\pi ijx/L}$, so the physical frequency operator is
$(2\pi/L)D$ and $\eta^*f=\sqrt L f(0)$ for a truncated Fourier vector.
\end{definition}

\begin{definition}[One-sided window distribution and correlation form]
\label{def:ccm-tn-window-form}
Let $f,g$ be finite sums of the modes $u_j$ of
\cref{def:ccm-tn-loewner}, extended by zero outside $[0,L]$.
Write $f^*(x)=\overline{f(-x)}$ and
$(f^**g)(y)=\int\overline{f(x)}g(x+y)\,dx$.
For a real distribution $\mathcal D$ on the closed half interval,
acting on $C^\infty([0,L])$, define
\[
 Q(f,g)=\mathcal D\bigl((f^**g)(y)+(f^**g)(-y)\bigr),\qquad
 (Q_N)_{jk}=Q(u_j,u_k).
\]
The one-sided derivatives of the symmetrised correlation at $0$ and
$L$ are part of this prescription. A globally even distribution on
smooth full-line tests does not by itself specify these endpoint data.
\end{definition}

\begin{definition}[Continuous features and three meanings of kernel]
\label{def:ccm-tn-continuous}
For a finite-dimensional unitary group $U(t)=e^{-itH}$, $H=H^*$, on a retained space, put
$F_f=\int f(t)U(t)\,dt$ for a compactly supported integrable test function.
For a positive spectral measure $\TNmeasure$, the spectral feature is
$\widehat f(\omega)=\int f(t)e^{-i\omega t}\,dt$, when this belongs
to $L^2(\TNmeasure)$. These are separate constructions until a spectral
representation identifies them.
The \emph{distribution kernel} of a Weil form is the distribution used
to pair two test functions through their correlation; its
\emph{nullspace} is the radical of that form. The \emph{Dirichlet kernel}
$\sum_{j=-N}^N e^{2\pi ijx/L}$ approximates evaluation at the boundary.
The archimedean gamma kernels are constituents of the input distribution,
not any of these nullspaces. Endpoint distributions retain the one-sided
prescription used to define the window form.
\end{definition}
